% Fragmento integrable; no documento autónomo y no orden de compilación.
% Fecha: 2026-09-30. No modifica fuentes publicadas.
% Procedencia: reunión de RETROACCION_PCH_Y_RESTRICCIONES_VARIACIONALES.md,
% PRECUANTIZACION_PCH_Y_FIDELIDAD.md y CIERRE_CANONICO_TOTAL_Y_REPRESENTACION.md.
% Prerrequisitos de inserción: núcleo común APP--TRIT--TPK; VIII/30d,
% VIII/31 y VIII/33; lectores de acción, gamma y acoplamientos.
% Los nuevos labels usan exclusivamente el prefijo hmtcuatro:.
\section{Acción conjunta, retroacción y cierre canónico de las restricciones}
\label{hmtcuatro:accion-retroaccion-restricciones}

La secuencia causal de este desarrollo es
\[
 \mathrm{APP}\longrightarrow\mathrm{TRIT}\longrightarrow\mathrm{TPK}
 \longrightarrow\text{estado enriquecido}
 \longrightarrow\text{estructura discreta conjunta del continuo}.
\]
Se conservan las hojas, la orientación, los residuos, el acarreo, las rutas,
la frontera y la memoria de esa construcción. La escala de acción y los
acoplamientos son salidas HMT recibidas antes de emplear el lenguaje
variacional y simpléctico. No se seleccionan por valores objetivo ni por
la conclusión de curvatura que se va a demostrar. El retorno de fase
tampoco reinicia el estado enriquecido.

La realización geométrica previamente construida proporciona una
cotetrada no degenerada \(e\), su conexión métrica y el levantamiento
espinorial del transporte. Aquí esa geometría se varía junto con los
campos materiales. Reuniremos en una sola demostración la acción común,
su reducción torsional, la transformación de Legendre, la primera clase,
la propagación de las restricciones y la planitud característica de la
conexión canónica que resulta de la misma acción. Esta última realización
no se identifica por su nombre con otro Hamiltoniano interactuante.

\subsection{Estructura de partida y acción con fuente completa}
\label{hmtcuatro:datos-accion-total}

Sea \(M\) una región lisa orientada de dimensión cuatro con estructura
de espín, firma \((-+++)\) y una foliación espacial regular. La cotetrada
toma valores en la recta de longitud de la realización recibida.
Escribimos
\[
 \Sigma^{IJ}=e^I\wedge e^J,\qquad
 F_\omega=\mathrm d\omega+\omega\wedge\omega,\qquad
 \mathsf P_\gamma=\star+\gamma^{-1}I,\qquad \star^2=-I.
\]
La estrella de esta fórmula actúa sobre bivectores internos; no es la
dualidad exterior de formas. Durante la variación permanecen fijas las
secciones recibidas \(c>0,\hbar>0,G>0,\gamma\in\mathbb R\setminus\{0\}\)
y \(\Lambda\). Con \(\kappa=8\pi G/c^4\), la acción total es
\begin{equation}
 \begin{split}
 S_{\rm tot}={}&\frac1{2\kappa c}
   \int_M\Sigma_{IJ}\wedge(\mathsf P_\gamma F_\omega)^{IJ}
 -\frac{\Lambda}{\kappa c}\int_M\operatorname{vol}_e\\
 &+S_{\rm YM}[e,A]+S_H[e,A,H]
   +S_W[e,\omega,A,\Psi]+S_Y[e,H,\Psi]+S_{\partial}.
 \end{split}
 \label{hmtcuatro:accion-total}
\end{equation}

\begin{definition}[Unificación variacional y canónica]
\label{hmtcuatro:def-unificacion}
La unificación variacional y canónica HMT es la composición de las realizaciones geométrica, cromática y electrodébil del mismo estado enriquecido en la acción común \eqref{hmtcuatro:accion-total}, cuyas variaciones producen fuentes conjuntas y cuya reducción conserva una misma álgebra de restricciones y su propagación. La componente electromagnética pertenece a la realización electrodébil. El teorema~\ref{hmtcuatro:teorema-canonico-total} establece esta composición en el dominio regular y con las condiciones de borde declaradas.
\end{definition}
Los acoplamientos electrodébiles $g,g'$ y el autoacoplamiento $\lambda_H$ se reciben de la construcción de \ref{amp-higgs-seccion} y del teorema~\ref{amp-higgs-inversa-teorema}; la variación presente mantiene fijas esas secciones. Su comparación entre secciones es una operación distinta, conservada en el desarrollo del horizonte de compatibilidad.


El símbolo \(H\) sin subíndice designa aquí el campo Higgs, no un
Hamiltoniano. Todas las contracciones espaciotemporales de la materia
utilizan la misma \(e\). El problema de borde es el recibido de la
variación Cartan--Holst: variaciones interiores, datos de borde fijos
compatibles o una completación que anule el mismo término fronterizo.

Para fijar el contenido de \eqref{hmtcuatro:accion-total}, las conexiones
internas actúan sobre el fibrado quiral
\[
 (S_L\otimes E_L)\oplus(S_R\otimes E_R).
\]
No se sustituye ese fibrado por un acoplamiento vectorial. En la
normalización entera \(y=6Y\), los multipletes recibidos tienen pesos
\(1,4,-2,-3,-6\), respectivamente para \(q_L,u_R,d_R,\ell_L,e_R\);
el doblete \(H\) tiene peso \(3\). El color actúa en el factor
\(\mathbb C^3\) de los quarks y trivialmente en los leptones. El factor
de familias y sus mapas \(Y_u,Y_d,Y_e\) se conservan. No se añade en
este bloque un compañero derecho a un multiplete que no lo tenga.

Con \(\widetilde H=i\sigma_2\overline H\), el mapa material es
\begin{equation}
 \begin{split}
 \mathsf Y_q(H)(u,d)&=
       \widetilde H\otimes Y_u u+H\otimes Y_d d,\\
 \mathsf Y_\ell(H)e&=H\otimes Y_e e .
 \end{split}
 \label{hmtcuatro:yukawa}
\end{equation}
Se entiende la identidad sobre el factor de color en la primera línea.
El término de acción es
\(-(\bar\Psi_L\mathsf Y(H)\Psi_R+\mathrm{h.c.})\), con las secciones
dimensionales recibidas. La equivariancia no es sólo nominal: para
\(U\in SU(2)\), \(i\sigma_2\overline U=U i\sigma_2\), de modo que
\(\widetilde H\mapsto U\widetilde H\). Los pesos de las tres columnas
son \(-3+4=1\), \(3-2=1\) y \(3-6=-3\). El color se entrelaza por la
identidad y los mapas de familia conmutan con la acción interna.
Por tanto, para cada transformación interna \(u\),
\begin{equation}
 \mathsf Y(uH)\rho_R(u)=\rho_L(u)\mathsf Y(H).
 \label{hmtcuatro:equivariancia-yukawa}
\end{equation}
Si la realización material precedente contiene además otros
multipletes, sus mapas se incorporan con su identidad de equivariancia
propia; esta prueba no inventa estados nuevos.

En una carta de unidades \(\hbar=c=1\), la densidad material que resume
estas acciones tiene la forma
\begin{equation}
 \begin{split}
 \mathcal L_m={}&-\frac14 F_{B,\mu\nu}F_B^{\mu\nu}
 -\frac14\sum_a F_{W,\mu\nu}^aF_W^{a,\mu\nu}
 -\frac14\sum_a F_{C,\mu\nu}^aF_C^{a,\mu\nu}\\
 &-(\nabla_\mu H)^\dagger\nabla^\mu H-V(H)
 +\mathcal L_W(\Psi_L,\nabla_L)+\mathcal L_W(\Psi_R,\nabla_R)\\
 &-\bigl(\bar\Psi_L\mathsf Y(H)\Psi_R+\mathrm{h.c.}\bigr).
 \end{split}
 \label{hmtcuatro:densidad-material}
\end{equation}
Las curvaturas incluyen sus acoplamientos; éstos y las normalizaciones
de componentes son los de sus propietarios. La escritura en esa carta
no elimina la sección de acción ni recalibra los coeficientes.
En particular, no se introduce aquí un valor numérico nuevo para el
acoplamiento fuerte. La conexión de espín de \(\nabla_L,\nabla_R\)
es la de \(e,\omega\); las conexiones \(B,W,C\) actúan en las
representaciones recién descritas. La constante de \(V\) se conserva,
pues su variación de volumen pertenece a la fuente gravitatoria.

La parte espinorial conserva las convenciones explícitas del desarrollo
precedente: \(\{g^I,g^J\}=2\eta^{IJ}I\),
\(\mathbb B_{IJ}=\frac14[g_I,g_J]\), \(A_D=-ig^0\),
\(\bar\psi=\psi^\dagger A_D\) y
\(D\psi=\mathrm d\psi+\frac12\omega^{IJ}\mathbb B_{IJ}\psi+A\psi\).
Su cinética simétrica es
\[
 \frac{\hbar}{2}
 [\bar\psi g^I D\psi-(D\bar\psi)g^I\psi]\wedge\eta_I,
 \qquad \eta_I=\iota_{E_I}\operatorname{vol}_e .
\]
La suma se toma sobre los campos realmente presentes, con sus
proyecciones quirales, no sobre una duplicación vectorial del catálogo.
Los términos Yukawa son algebraicos y no dependen de \(\omega\).
La corriente total queda fijada por
\begin{equation}
 \delta_\omega S_m=-\frac12\int_M\delta\omega^{IJ}\wedge\sigma_{IJ},
 \qquad
 \sigma_{JK}=\sum_s-\frac{\hbar}{2}
   \bar\psi_s\{g^I,\mathbb B_{JK}\}\psi_s\,\eta_I.
 \label{hmtcuatro:corriente-total}
\end{equation}

\subsection{Una demostración de compatibilidad, propagación y cierre}
\label{hmtcuatro:demostracion-unica}

\begin{theorem}[Compatibilidad canónica de la acción conjunta]
\label{hmtcuatro:teorema-canonico-total}
Considérese \eqref{hmtcuatro:accion-total} en el dominio no degenerado
anterior, con materia afín en la contorsión, y en una carta regular de
la transformación de Legendre restringida. Reténganse las restricciones
fermiónicas de primer orden y elimínense las de segunda clase mediante
su corchete de Dirac. Para parámetros de soporte interior, o con el
borde compatible sin cargos impropios, se cumplen conjuntamente:
la ecuación métrica tiene la fuente material total de esa misma acción;
sus generadores de difeomorfías y calibre son de primera clase; las
restricciones se propagan sobre toda solución acoplada regular; y el
potencial canónico total induce una conexión precuántica plana en las
direcciones características de la superficie regular de restricciones.

La afirmación es sobre esta acción y esta realización canónica.
No identifica automáticamente su conexión con un Hamiltoniano
interactuante distinto ni prueba su límite espacial--quiral.
\end{theorem}

\begin{proof}
\par\medskip\noindent\textbf{1. Eliminación torsional sin duplicación.}\quad
Escribamos \(\omega=\mathring\omega(e)+\mathfrak k\).
La expansión \(F_\omega=F_{\mathring\omega}
+D_{\mathring\omega}\mathfrak k+\mathfrak k\wedge\mathfrak k\)
y \(D_{\mathring\omega}\Sigma=0\) separan la contribución de borde
\[
 S_{\partial,\mathfrak k}
 =\frac1{2\kappa c}\int_{\partial M}
                  \Sigma_{IJ}\wedge(\mathsf P_\gamma\mathfrak k)^{IJ}
\]
de la densidad algebraica
\begin{equation}
 \mathcal Q_{e,\gamma}(\mathfrak k)
       -\frac12\mathfrak k^{IJ}\wedge\sigma_{IJ},\qquad
 \mathcal Q_{e,\gamma}(\mathfrak k)=\frac1{2\kappa c}
       \Sigma_{IJ}\wedge
       \mathsf P_\gamma(\mathfrak k\wedge\mathfrak k)^{IJ}.
 \label{hmtcuatro:cuadratica-contorsion}
\end{equation}
La variación original da
\(D_\omega(\mathsf P_\gamma\Sigma)=\kappa c\,\sigma\).
La inversa de Holst y la inversa torsión--contorsión del desarrollo
precedente son invertibles para el coframe y \(\gamma\) declarados.
Así determinan una única \(\mathfrak k_*[e,\sigma]\), lineal en
\(\sigma\). La variación radial de
\eqref{hmtcuatro:cuadratica-contorsion} en su estacionario da
\(2\mathcal Q(\mathfrak k_*)=\frac12\mathfrak k_*^{IJ}\wedge\sigma_{IJ}\).
En consecuencia,
\begin{equation}
 S_{\rm red}=\frac1{2\kappa c}\int_M(R-2\Lambda)\operatorname{vol}_e
       +S_m^{\rm LC}
       -\frac14\int_M\mathfrak k_*^{IJ}\wedge\sigma_{IJ}
       +S_{\partial,\rm red}.
 \label{hmtcuatro:accion-reducida}
\end{equation}
La primera identidad de Bianchi anula el término Holst de
Levi--Civita en esta escritura. La regla de la cadena conserva las
ecuaciones de \(e,A,H,\Psi\), pues el coeficiente de
\(\delta\mathfrak k_*\) se anula por estacionariedad.
Como \(\sigma=\sum_s\sigma_s\), la contribución efectiva incluye
\(\mathfrak k_*[\sigma_s]\wedge\sigma_t\) para \(s\ne t\).
No se reemplaza por cuadrados de especies aisladas ni se conserva
simultáneamente la misma \(\mathfrak k\) como variable independiente.

\par\medskip\noindent\textbf{2. Fuente completa y retroacción.}\quad
Definamos el tensor de acción por la variación de todos los términos
materiales de \eqref{hmtcuatro:accion-reducida}:
\begin{equation}
 \delta_g S_{m,\rm red}=-\frac12\int_M
       \mathcal T^S_{\mu\nu}\,\delta g^{\mu\nu}\operatorname{vol}_e,
 \qquad
 E_{\mu\nu}:=G_{\mu\nu}+\Lambda g_{\mu\nu}
                   -\kappa c\,\mathcal T^S_{\mu\nu}.
 \label{hmtcuatro:fuente-completa}
\end{equation}
La ecuación métrica es \(E_{\mu\nu}=0\).
El factor \(\kappa c\) corresponde a la cotetrada de longitud;
si \(T^{\rm phys}=c\mathcal T^S\), se escribe \(\kappa T^{\rm phys}\).
Se varían también los campos materiales y las conexiones internas.
Por ello los campos determinan la fuente y la geometría que resuelve
la ecuación determina a su vez sus conexiones, volumen y evolución.
Ésta es retroacción, no evolución sobre un fondo fijado.
La variación incluye el volumen del potencial Higgs y la dependencia
métrica de la corriente torsional, además de sus términos cruzados.

\par\medskip\noindent\textbf{3. Legendre restringida y reducción canónica.}\quad
El potencial canónico total se obtiene de
\(\delta L=\mathcal E\delta z+\mathrm d\theta(\delta z)\),
integrando \(\theta\) sobre la hoja espacial y efectuando la reducción
de segunda clase. Su parte geométrica antes de esa reducción es
\[
 \theta_{\rm geom}(\delta)=
   \frac1{2\kappa c}(\mathsf P_\gamma\Sigma)_{IJ}
                       \wedge\delta\omega^{IJ}.
\]
La descomposición \(F_{0a}=\dot\omega_a-D_a\omega_0\), junto con la
misma descomposición de las acciones materiales, produce
\begin{equation}
 S_{\rm can}=\int d\tau\,
 \bigl[\Theta_{\rm red}(\dot z)-\mathcal C[N]-\mathcal D[\xi]
       -\mathcal G_{\rm L}[\lambda]-\mathcal G_{\rm int}[\chi]\bigr]
       +S_{\partial,\rm red}.
 \label{hmtcuatro:legendre-total}
\end{equation}
No se invierten artificialmente las velocidades ausentes de lapse,
shift o conexiones temporales.

La contorsión algebraica tiene restricciones
\(\pi_{\mathfrak k}=0\) y
\(\partial\mathcal Q/\partial\mathfrak k-\sigma/2=0\).
Su matriz de corchetes tiene el bloque
\[
 \begin{pmatrix}0&-M\\ M^T&B\end{pmatrix},
 \qquad M=\frac{\partial^2\mathcal Q}{\partial\mathfrak k^2}.
\]
La inversa torsional hace invertible \(M\), y por tanto esta matriz.
Son restricciones de segunda clase. Se usa
\begin{equation}
 \{F,G\}_D=\{F,G\}
       -\{F,\chi_a\}(C^{-1})^{ab}\{\chi_b,G\}.
 \label{hmtcuatro:corchete-dirac}
\end{equation}
En las variables que no dependen de la pareja contorsión--momento,
el bloque inferior derecho de \(C^{-1}\) es cero. No aparece por esa
eliminación un corchete suplementario entre ellas. Se retienen,
separadamente, las restricciones espinoriales de primer orden; en la
fase graduada se emplea el corchete graduado y los generadores pares.
\(\Theta_{\rm red}\) designa el resultado completo, no sólo su parte
geométrica.

La inversión de Legendre en direcciones regulares recupera la misma
acción. Por ejemplo, \(N[p^2/(2w)+wV+H_{\rm grad}]\), con
\(w=\sqrt{\det q}\), da \(p=(w/N)D_\tau\phi\), y por sustitución
\[
 L_H=\frac{w}{2N}|D_\tau\phi|^2-NwV-NH_{\rm grad}.
\]
\(D_\tau\) conserva shift y conexión temporal; \(w\) no se congela.
Para \(K_{ij}=(\dot q_{ij}-\mathcal L_\xi q_{ij})/(2N)\), la
contribución gravitatoria al momento y su inversa son
\begin{equation}
 \Pi^{ij}=\frac{\sqrt q}{2\kappa c}(K^{ij}-q^{ij}K),\qquad
 K_{ij}=\frac{2\kappa c}{\sqrt q}
                  (\Pi_{ij}-\tfrac12q_{ij}\Pi).
 \label{hmtcuatro:legendre-geometrica}
\end{equation}
De ellas resulta
\[
 \mathcal C_{\rm grav}
 =\frac{2\kappa c}{\sqrt q}
       (\Pi^{ij}\Pi_{ij}-\tfrac12\Pi^2)
  -\frac{\sqrt q}{2\kappa c}(R^{(3)}-2\Lambda).
\]
Si la carta de espinores aporta un desplazamiento material del momento
total \(p\), se usa \(\Pi=p-p_m\); no se descarta \(p_m\).
El carácter indefinido de la forma de DeWitt no impide esta inversa algebraica.
No se exige una inversa de velocidades a la cinética fermiónica de
primer orden.

\par\medskip\noindent\textbf{4. Identidad Noether--Legendre y primera clase.}\quad
La coincidencia de los generadores con las restricciones de la misma
acción puede verse antes de imponer sus ecuaciones.
Pónganse \(a=(2\kappa c)^{-1}\), \(B=\mathsf P_\gamma\Sigma\),
\(u^I=\iota_\zeta e^I\) y \(\lambda^{IJ}=\iota_\zeta\omega^{IJ}\).
De \(\mathcal L_\zeta\omega=\iota_\zeta F+D_\omega\lambda\) se obtiene
\begin{equation}
 \begin{split}
 j_\zeta^{\rm geom}
 ={}&\mathrm d(a B_{IJ}\lambda^{IJ})
   -a(D_\omega B_{IJ})\lambda^{IJ}\\
 &-2a u^Ie^J\wedge(\mathsf P_\gamma F)_{IJ}
       +2a\Lambda u^I\eta_I .
 \end{split}
 \label{hmtcuatro:noether-legendre}
\end{equation}
La parte material añade exactamente sus contribuciones a las
ecuaciones de \(e,\omega,A\), sus generadores verticales y su borde.
Descomponer la misma expresión en una hoja espacial da
\eqref{hmtcuatro:legendre-total}. La reducción estacionaria y la
reducción de Dirac transportan esas transformaciones, en vez de
sustituirlas por momentos de una acción externa.

Con la convención canónica usual para las restricciones, los corchetes
totales son, módulo los generadores Gauss conservados,
\begin{equation}
 \begin{aligned}
 \{\mathcal D[\xi],\mathcal D[\eta]\}_D
    &\simeq\mathcal D[[\xi,\eta]],\\
 \{\mathcal D[\xi],\mathcal C[N]\}_D
    &\simeq\mathcal C[\mathcal L_\xi N],\\
 \{\mathcal C[N],\mathcal C[M]\}_D
    &\simeq\mathcal D[
       q^{ij}(N\partial_jM-M\partial_jN)\partial_i].
 \end{aligned}
 \label{hmtcuatro:algebra-restricciones}
\end{equation}
En efecto, las transformaciones tangenciales tienen el corchete de
Lie, y transportan el lapse como escalar. Para las normales,
diferenciar su ortogonalidad con los vectores tangentes y
\(g(n,n)=-1\) determina la variación tangencial de \(n\);
con \(K_{ij}\) de \eqref{hmtcuatro:legendre-geometrica} se obtiene
\(\delta_N n=\operatorname{grad}_qN\). El conmutador de deformaciones,
convertido al corchete canónico con estas convenciones, da la última
línea de \eqref{hmtcuatro:algebra-restricciones}. El uso de transporte
covariantizado conserva además las rotaciones de marco y de calibre.
Las transformaciones son tangentes a la superficie de restricciones
de la acción invariante, de donde resulta la primera clase.
Las funciones \(q^{ij}\) dependen del estado y no se congelan.
Con soporte interior no se añade un cargo central de borde; un
generador impropio con cargo no nulo no se declara restricción cero.

\par\medskip\noindent\textbf{5. Propagación en la geometría que responde a la materia.}\quad
La invariancia por difeomorfías de \(S_{m,\rm red}\), evaluada sobre
sus ecuaciones materiales, da por integración por partes
\(\mathring\nabla^\mu\mathcal T^S_{\mu\nu}=0\). La identidad usa el
tensor completo de \eqref{hmtcuatro:fuente-completa}; no postula
conservación separada de cada contribución.
Bianchi implica entonces
\(\mathring\nabla_\mu E^{\mu\nu}=0\), aun antes de imponer
todas las ecuaciones métricas.

En una carta gaussiana local,
\(ds^2=-d\tau^2+q_{ij}dx^idx^j\), impongamos las ecuaciones
evolutivas \(E_{ij}=0\) y pongamos
\(C=E_{00}\), \(C_i=E_{0i}\),
\(K_{ij}=\dot q_{ij}/2\), \(K=q^{ij}K_{ij}\).
Aquí \(K\) es sólo la traza de la segunda forma fundamental.
Las dos componentes de la identidad divergente son
\begin{equation}
 \partial_\tau C-D_iC^i+KC=0,\qquad
 \partial_\tau C_i+KC_i=0.
 \label{hmtcuatro:propagacion}
\end{equation}
Para comprobar los signos se usan
\(E^{00}=C\), \(E^{0i}=-C^i\), \(E^{ij}=0\),
\(\Gamma^i{}_{0j}=K^i{}_j\) y
\(\Gamma^\mu{}_{\mu0}=K\). Al bajar el índice de la ecuación
espacial, \(\dot q_{ij}=2K_{ij}\) cancela los dos términos de conexión.
Si inicialmente \(C_i=0\), su ecuación homogénea mantiene \(C_i=0\);
la primera queda \(\dot C+KC=0\) y conserva \(C=0\).
\(q\) y \(K\) son los de la solución acoplada. La elección gaussiana
prueba la afirmación local; su forma tensorial permite cambiar lapse
y shift. Análogamente, la identidad gauge de Noether relaciona la
divergencia covariante de la ecuación de calibre con las ecuaciones
materiales y propaga Gauss cuando se cumplen las ecuaciones espaciales.
No se infiere existencia global ni ausencia de singularidades de
toda solución.

\par\medskip\noindent\textbf{6. La misma acción induce el cierre característico.}\quad
Sea \(\mathcal P_{\rm red}\) la carta canónica regular obtenida y
\(\mathcal Z=\{C_A=0\}\) su superficie regular de restricciones
totales, incluida Gauss. Usamos un símbolo distinto de \(\Sigma^{IJ}\)
para esa superficie. Se conserva el potencial \(\Theta=\Theta_{\rm red}\)
de \eqref{hmtcuatro:legendre-total}, incluidos los términos materiales
y mixtos. Fijamos
\begin{equation}
 \Omega=-\mathrm d_{\mathcal P}\Theta,\qquad
 \iota_{X_f}\Omega=\mathrm d_{\mathcal P}f,\qquad
 \{f,g\}_D=\Omega(X_f,X_g).
 \label{hmtcuatro:convenciones-simpecticas}
\end{equation}
En una pareja canónica, \(\Theta=p\,\mathrm dq\) da
\(X_f=f_p\partial_q-f_q\partial_p\), \(\{q,p\}_D=1\),
\(X_f(g)=-\{f,g\}_D\) y
\([X_f,X_g]=-X_{\{f,g\}_D}\). En la fase graduada esta escritura
se aplica a las funciones y restricciones pares pertinentes.

En la línea de fase trivializada sobre esta carta, construimos
\begin{equation}
 \nabla_X=X-\frac{i}{\hbar}\Theta(X)
       =\delta_X+\frac{i}{\hbar}H_X^{\rm can},
 \qquad H_X^{\rm can}=-\Theta(X).
 \label{hmtcuatro:conexion-canonica}
\end{equation}
No se escoge una conjugación arbitraria para hacerla plana:
su coeficiente proviene del potencial de la acción.
Para cualesquiera campos \(X,Y\) de la carta, el cálculo da
\begin{equation}
 \begin{aligned}
 \mathcal F^{\rm can}_{X,Y}
 &=XH_Y^{\rm can}-YH_X^{\rm can}-H_{[X,Y]}^{\rm can}
       +\frac{i}{\hbar}[H_X^{\rm can},H_Y^{\rm can}]\\
 &=-X\Theta(Y)+Y\Theta(X)+\Theta([X,Y])\\
 &=-\mathrm d_{\mathcal P}\Theta(X,Y)=\Omega(X,Y).
 \end{aligned}
 \label{hmtcuatro:curvatura-calculada}
\end{equation}
Los \(H_X^{\rm can}\) son multiplicadores de la línea, por lo que
su conmutador es cero; las derivadas de los coeficientes y de los
campos no se suprimen.

La primera clase ya probada significa
\(\{C_A,C_B\}_D=f_{AB}{}^D C_D\).
Para \(V\) tangente a \(\mathcal Z\),
\(\Omega(X_{C_A},V)=\mathrm d_{\mathcal P}C_A(V)=0\).
Los campos de las restricciones son, por tanto, característicos.
Bajo la regularidad declarada generan la distribución característica
de la superficie coisótropa. La tensorialidad de \(\Omega\) extiende
el cálculo a todas sus combinaciones suaves y prueba
\begin{equation}
 \left.\mathcal F^{\rm can}_{X,Y}\right|_{\mathcal Z}=0
 \quad\hbox{para }X,Y\hbox{ característicos}.
 \label{hmtcuatro:cierre-caracteristico}
\end{equation}
Para dos lapsos se incluye en \([X,Y]\) la deformación tangencial
y las acciones Gauss de \eqref{hmtcuatro:algebra-restricciones}.
Si se pasa al cociente gauge, la igualdad se expresa módulo esas
acciones verticales. No se afirma que \(\Omega\) sea cero en toda
la fase ni que toda holonomía global de una hoja sea trivial.

\par\medskip\noindent\textbf{7. Expresión precuántica del mismo cierre.}\quad
La realización diferencial asociada a
\eqref{hmtcuatro:conexion-canonica}, sobre secciones suaves donde
están definidos los productos, es
\begin{equation}
 Q_{\rm pre}(f)=-i\hbar\nabla_{X_f}+f
            =-i\hbar X_f+f-\Theta(X_f).
 \label{hmtcuatro:operador-precuantico}
\end{equation}
Con \(R^\nabla=(i/\hbar)\Omega\), la expansión del conmutador y
\(X_f(g)=-\{f,g\}_D\) dan
\begin{equation}
 [Q_{\rm pre}(f),Q_{\rm pre}(g)]
           =i\hbar Q_{\rm pre}(\{f,g\}_D).
 \label{hmtcuatro:representacion-precuantica}
\end{equation}
La regla de producto que conserva las funciones de estructura es
\begin{equation}
 Q_{\rm pre}(aC)
   =aQ_{\rm pre}(C)+C\bigl(Q_{\rm pre}(a)-a\bigr).
 \label{hmtcuatro:funciones-estructura}
\end{equation}
Se deduce de \(X_{aC}=aX_C+CX_a\) y de
\eqref{hmtcuatro:operador-precuantico}. No se elimina el segundo
término fuera de \(\mathcal Z\).
Sobre \(\mathcal Z\), \(Q_{\rm pre}(C_A)=-i\hbar\nabla_{X_{C_A}}\);
así \eqref{hmtcuatro:representacion-precuantica} y
\eqref{hmtcuatro:cierre-caracteristico} son las expresiones operatoria
y geométrica del mismo cierre canónico, no dos incorporaciones
independientes de gravedad.
\end{proof}

\subsection{Dominio del resultado y comparación de realizaciones}
\label{hmtcuatro:alcance-canonico}

La prueba reúne una acción con gravedad dinámica y fuente conjunta,
su Legendre y su conexión característica. El calificativo
\emph{arbitrario} para las deformaciones se refiere a los lapsos y
desplazamientos admisibles que generan direcciones características en
esa región regular; no a cualquier vector de fase, cualquier borde
o cualquier solución singular. La planitud local conserva la posible
monodromía global y no identifica dos historias TPK por tener iguales
extremos.

El coeficiente \(H_X^{\rm can}\) de una conexión de línea no es por
definición el operador de reloj y respuesta \(H_{\rm clk}+D^*WD\),
ni su extensión interactuante. \(Q_{\rm pre}(C)\) incluye además
una derivación en la fase; tampoco es un nombre alternativo de esos
operadores. Para identificar una realización concreta mediante un
mapa \(I\) se requiere la igualdad tipada
\[
 \nabla_N^{\rm op}I=I\nabla_N^{\rm can}
\]
sobre sus dominios efectivos. Si \(I\) depende de la geometría, su
derivada forma parte de esta igualdad.

La distinción tiene un control exacto: en la carta
\(\Theta=p\,\mathrm dq\),
\begin{equation}
 Q_{\rm pre}(p^2/2)=-i\hbar p\,\partial_q-p^2/2.
 \label{hmtcuatro:control-polarizacion}
\end{equation}
Al actuar sobre la función \(1\), produce \(-p^2/2\).
Por tanto no preserva por mera restricción la polarización de
funciones independientes de \(p\), y no es allí
\(-\hbar^2\partial_q^2/2\). El control no refuta el Hamiltoniano
recibido: impide sustituir una representación por otra sin su mapa.

No se ha supuesto una medida de Liouville infinito-dimensional ni
se deduce autoadjunción de las identidades diferenciales anteriores.
Si inclusiones de refinamiento entrelazan las conexiones y conservan
un núcleo común, sus curvaturas se entrelazan por composición; un
límite exige además los controles de dominio y convergencia propios.
No se presenta aquí esa condición como comprobada para el operador
interactuante y el límite espacial--quiral conjunto. Las construcciones
cuánticas anteriores se conservan; el resultado presente precisa
exactamente el cierre canónico que se acaba de demostrar.
